\chapter{Présentation des modèles}
\label{chap:presentation_des_modeles}

%%%%%%%%%%%%%%%%%%%%%%%%%%%%%%%%%%
% Présentation des modèles de taux
%%%%%%%%%%%%%%%%%%%%%%%%%%%%%%%%%%
\section{Présentation des modèles de taux}
% Modèles à facteur court
% Modèles de structure par terme
% Modèles de taux stochastiques multi-facteurs

Dans cette partie, nous allons présenter les différents modèles qui peuvent être utilisés pour modéliser l'évolution des taux d'intérêt. Contrairement aux actions, les taux ne sont pas décrits par une seule variable observable, mais par une courbe des taux dépendant de l'échéance. Cette spécificité a conduit au développement de plusieurs familles de modèles parmi lesquelles on retrouve les modèles de taux court, les modèles multifactoriels et le modèle de Heath-Jarrow-Morton (HJM). 

\subsection{Les modèles de taux court}

Les modèles de taux court, ou short-rate models, reposent sur l'idée que toute la dynamique de la courbe des taux peut être expliquée à partir de l'évolution d'un seul taux instantané de court terme, noté en général $r_t$. Ce taux court est supposé suivre un processus stochastique, et les prix des obligations zéro-coupon sont ensuite obtenus par actualisation sous une probabilité risque-neutre.

La forme générale d'un modèle à facteur court s'écrit :
\begin{equation}
dr_t = \mu(r_t, t) dt + \sigma(r_t, t) dW_t
\end{equation}
où $\mu$ est la dérive, $\sigma$ la volatilité, et $W_t$ un mouvement brownien standard.

Parmi les modèles à facteur court les plus connus, on peut citer :
\begin{itemize}
    \item Le modèle de Vasicek (1977) : modèle à retour à la moyenne avec une dynamique gaussienne.
    \item Le modèle de Cox-Ingersoll-Ross (CIR) : modèle proche de Vasicek mais avec une volatilité proportionnelle à la racine carrée du taux, ce qui garantit que les taux restent positifs.
    \item Le modèle de Hull-White (1990) : extension du modèle de Vasicek avec des paramètres dépendant du temps, permettant de mieux ajuster la courbe des taux initiale.
\end{itemize}

Ces modèles ont plusieurs avantages. Ils sont généralement simples et bien adaptés à la valorisation de nombreux produits dérivés de taux d'intérêt. En particulier, certains modèles permettent d'obtenir des formules fermées pour les prix des obligations zéro-coupon, ce qui facilite les calculs.
Cependant, leur principale limite tient au fait  qu'un seul facteur est souvent insuffisant pour reproduire fidèlement les mouvements de la courbe des taux observés sur les marchés. En effet, la courbe des taux se déforme selon plusieurs dimensions (niveau, pente, courbure), et un modèle à facteur court ne peut pas capturer ces différentes sources de variation de manière satisfaisante.

\subsection{Les modèles multifactoriels}

Les modèles monofactoriels sont souvent très vite limités lorsque l'on s'intéresse à des produits de taux faisant intervenir plusieurs maturités. En effet, les variations de taux zéro-coupon de différentes maturités ne sont pas parfaitement corrélées et la corrélation diminue très fortement plus les maturités sont éloignées. 

Afin de pallier cette limite, les premiers modèles multifactoriels ont proposé d'ajouter un taux long terme $l(t)$ au modèle de taux court. Ce type de modèle bifactoriel a été développé notamment par Brennan et Schwartz (1982) et par Schaefer et Schwartz (1984). Cependant, ces modèles on été abandonnés pour plusieurs raisons. Tout d'abord, ils ne permettent pas d'obtenir d'évaluation simple des obligations zéro-coupon, sauf au prix d'approximations analytiques. De plus, ils sont en contradiction avec un résultat de 1996, montrant qu'en l'absence d'opportunité d'arbitrage, le taux long terme ne peut pas diminuer.

Ainsi, les modèles factoriels utilisés aujourd'hui ont plutôt recours à des dynamiques de taux court à plusieurs facteurs stochastiques. Le modèle le plus utilisé dans cette catégorie est le modèle gaussien à deux facteurs présenté par Brigo et Mercurio (2005). Il est aussi couramment appelé G2++. 
On suppose que sous la mesure risque neutre, le taux court $r_t$ s'écrit :
\begin{equation}
r_t = x(t) + y(t) + \phi(t),    r(0) = r_0
\end{equation}
où les deux facteurs $x(t)$ et $y(t)$ sont solutions des équations différentielles stochastiques suivantes :
\begin{equation}
\begin{cases}
dx(t) = -a x(t) dt + \sigma dW_1(t),    x(0) = x_0 \\
dy(t) = -b y(t) dt + \eta dW_2(t),    y(0) = y_0
\end{cases}
\end{equation}
avec $a$, $b$, $\sigma$ et $\eta$ des paramètres positifs, $W_1(t)$ et $W_2(t)$ deux mouvements browniens corrélés de coefficient de corrélation $\rho$, et $\phi(t)$ une fonction déterministe permettant d'ajuster la courbe des taux initiale.

Ainsi, les modèles multifactoriels, et en particulier le modèle G2++, peuvent être vus comme un prolongement utile des approches monofactorielles. Cette approche permet généralement de mieux rendre compte des variations de la courbe des taux. 
Toutefois, leur construction à partir du taux court peut encore présenter certaines limites, notamment lorsque l'on cherche à décrire de manière directe la dynamique de l'ensemble de la structure par terme. Dans cette perspective, il peut être pertinent de s'intéresser au cadre de Heath-Jarrow-Morton (HJM), qui propose une modélisation fondée directement sur les taux forwards.

\subsection{Le modèle de Heath-Jarrow-Morton (HJM)}

Le modèle de Heath, Jarrow et Morton (HJM), introduit en 1992, propose un cadre de travail plus large que les modèles précédents. Son idée centrale est de prendre comme donnée la courbe initiale des taux forwards instantanés $f(0,T)$, puis de modéliser directement l'évolution de la courbe $f(t,T)$ pour l'ensemble des maturités $T$, plutôt que de partir uniquement du taux court. 

Pour chaque échéance T, le taux forward est donné par : 
\begin{equation}
    f(t,T) = f(0,T) + \int_0^t \alpha(s,T) ds + \int_0^t \sigma(s,T) dW_s \; avec \; 0\leq t \leq T
\end{equation}
Et sous la mesure risque neutre, la dynamique de $f(t,T)$ est donnée par l'équation différentielle suivante :
\begin{equation}
\begin{cases}
df(t,T) = \alpha(t,T) dt + \sigma(t,T) dW_t \\
f(0,T) = f^M(0,T)
\end{cases}
\end{equation}
où $\alpha(t,T)$ est la dérive, $\sigma(t,T)$ la volatilité, $W_t$ un mouvement brownien standard et $f^M(0,T)$ la courbe des taux forwards de maturité $T$ observée sur le marché.





%%%%%%%%%%%%%%%%%%%%%%%%%%%%%%%%%%
% Le Générateur de Scénario Economique (GSE)
%%%%%%%%%%%%%%%%%%%%%%%%%%%%%%%%%%
\section{Le Générateur de Scénarios Économiques (GSE)}

\subsection{Présentation générale du GSE}

Dans le cadre de Solvabilité II, les assureurs sont tenus d'évaluer leurs engagements en juste valeur, notamment au travers du calcul du \textit{Best Estimate}, qui repose sur la projection des flux futurs puis sur leur actualisation à partir de la courbe des taux sans risque. En assurance vie, cet exercice est particulièrement complexe, dans la mesure où les contrats comportent fréquemment des garanties financières ainsi que des options implicites dont la valeur dépend étroitement des conditions de marché. Or, cette valeur ne peut généralement pas être observée directement. Dans ce contexte, le \textbf{Générateur de Scénarios Économiques} (GSE) constitue un outil indispensable, puisqu'il permet de simuler de manière cohérente l'évolution future des principales variables financières, telles que les taux d'intérêt, les actions, les spreads de crédit ou encore l'inflation.

De manière opérationnelle, le GSE produit un grand nombre de scénarios économiques futurs, qui alimentent ensuite le modèle actif-passif (\textit{Asset-Liability Management}, ALM) utilisé par l'assureur afin de projeter les flux de trésorerie associés aux contrats. L'agrégation des résultats issus de l'ensemble de ces scénarios permet alors d'obtenir une estimation de la valeur moyenne des engagements. Le GSE occupe ainsi une place centrale, à la fois pour satisfaire les exigences prudentielles et pour appréhender de manière réaliste la valorisation des passifs en assurance vie.

Plus généralement, un Générateur de Scénarios Économiques peut être défini comme un dispositif de modélisation stochastique destiné à représenter l'évolution conjointe des principaux facteurs financiers et macroéconomiques. Il ne s'agit pas d'un modèle unique, mais d'un ensemble structuré de modèles articulés entre eux afin de produire un large éventail de trajectoires futures cohérentes. En pratique, trois grandes composantes sont généralement mobilisées : un modèle de taux d'intérêt, un modèle de rendement des actifs (action et immobilier) et un modèle d'inflation. En assurance vie, le recours à un GSE est d'autant plus important que la valeur des engagements dépend fortement de variables de marché de long terme, ainsi que de garanties optionnelles particulièrement sensibles à leur évolution.


\subsubsection{Le GSE risque neutre}

Le GSE risque neutre est défini sous une mesure de probabilité dite risque neutre, généralement notée $Q$. Son objectif n'est pas de représenter le scénario futur le plus vraisemblable d'un point de vue économique, mais de permettre la valorisation correcte de flux contingents à partir des informations incorporées dans les prix de marché. Dans ce cadre, les primes de risque sont neutralisées : les prix actualisés des actifs doivent être compatibles avec une hypothèse d'absence d'arbitrage, ce qui autorise la valorisation de produits complexes par le calcul d'une espérance actualisée. Sous cette mesure, le rendement espéré de l'ensemble des actifs est égal au taux sans risque.

En pratique, un GSE risque neutre est calibré à partir des données observées sur les marchés financiers, telles que la courbe des taux sans risque, les prix ou les volatilités implicites de swaptions, de caps, de floors ou encore d'options sur actions. Le calibrage vise à reproduire aussi fidèlement que possible les prix d'instruments liquides, afin de garantir une valorisation cohérente avec le marché (\textit{market consistent}) pour des engagements non cotés, comme les garanties intégrées aux contrats d'assurance vie. C'est dans cette perspective que ce type de GSE est mobilisé pour le calcul du \textit{Best Estimate}.

Un point méthodologique doit toutefois être souligné : un scénario généré sous la mesure risque neutre n'a pas vocation à être réaliste au sens historique du terme. Les trajectoires obtenues peuvent paraître peu plausibles lorsqu'elles sont comparées directement aux observations passées, précisément parce qu'elles intègrent les ajustements nécessaires pour retrouver les prix de marché. En d'autres termes, un GSE risque neutre est avant tout un outil de valorisation, et non un instrument de prévision économique.

\subsubsection{Le GSE monde réel}

Le GSE monde réel est construit sous la mesure de probabilité historique, généralement notée $P$. Contrairement au GSE risque neutre, son objectif n'est pas de valoriser des flux en cohérence avec les prix de marché, mais de représenter de manière plausible l'évolution future de l'environnement économique et financier. Il s'inscrit donc dans une logique de prévision économique et de mesure des risques. Dans ce cadre, l'enjeu est de produire des trajectoires qui reflètent au mieux les dynamiques susceptibles d'être observées à long terme, en tenant compte des caractéristiques empiriques des différentes variables modélisées.

La construction d'un GSE monde réel vise ainsi à reproduire un certain nombre de propriétés statistiques et économiques observées dans les données : niveau moyen des rendements, structure des volatilités, dépendances entre classes d'actifs, phénomènes de retour à la moyenne, asymétrie des distributions, épaisseur des queues ou encore effets de cycle. Elle peut également intégrer des hypothèses prospectives de long terme, par exemple sur la croissance, l'inflation, les taux d'intérêt ou les primes de risque. Le calibrage repose dès lors principalement sur des séries historiques, complétées, le cas échéant, par des hypothèses d'experts. À la différence du cadre risque neutre, il ne cherche donc pas prioritairement à reproduire les prix observés des instruments dérivés ni les volatilités implicites de marché.

Dans une perspective actuarielle, le GSE monde réel occupe une place centrale dans tous les exercices où l'on cherche à analyser la distribution future des résultats et non à obtenir une valeur de marché à une date donnée. Il est ainsi mobilisé dans les études d'\textit{Asset-Liability Management} (ALM), ainsi que dans le cadre de l'ORSA. Son usage est donc particulièrement pertinent dès lors que l'on s'intéresse à la trajectoire future de l'entreprise, à ses équilibres financiers de long terme et à sa capacité à respecter durablement ses contraintes prudentielles.

Sur le plan méthodologique, le GSE monde réel doit avant tout produire des scénarios crédibles du point de vue économique. La qualité d'un tel générateur se mesure à sa faculté à reproduire des comportements réalistes et cohérents entre variables. Cette exigence est essentielle en assurance vie, où les engagements s'inscrivent dans un horizon long et où les décisions de gestion dépendent fortement des interactions entre rendement des actifs, évolution des taux, inflation, comportement des assurés (rachats) et contraintes réglementaires.

\subsubsection{Les principaux points de différence entre les deux approches}

La différence fondamentale entre un GSE en monde réel et un GSE en risque neutre ne tient pas d'abord à la forme des trajectoires simulées, mais à la mesure de probabilité retenue, et donc à la manière dont les états futurs sont pondérés. Autrement dit, ce ne sont pas nécessairement les scénarios eux-mêmes qui diffèrent le plus, mais la probabilité qui leur est associée. Sous la mesure risque neutre, certains scénarios défavorables reçoivent un poids plus important, car les prix de marché reflètent l'aversion au risque des investisseurs. À l'inverse, sous la mesure monde réel, les probabilités cherchent à traduire la fréquence attendue ou la vraisemblance économique des événements.

Ces différences peuvent être synthétisées ainsi :
\begin{itemize}
    \item \textbf{Finalité :} le GSE en risque neutre est utilisé à des fins de valorisation, tandis que le GSE en monde réel sert à la projection économique et à l'évaluation des risques.
    \item \textbf{Mesure de probabilité :} le cadre risque neutre repose sur la mesure $Q$, cohérente avec les prix observés sur les marchés ; le cadre monde réel repose sur la mesure $P$, fondée sur une interprétation historique ou économique des probabilités.
    \item \textbf{Calibrage :} un GSE en risque neutre est calibré à partir des prix de marché et des volatilités implicites, alors qu'un GSE en monde réel s'appuie principalement sur les séries historiques et les relations économiques observées.
    \item \textbf{Contraintes de modélisation :} le risque neutre doit respecter la cohérence de marché, l'absence d'opportunités d'arbitrage et les propriétés de martingalité ; le monde réel vise avant tout à reproduire des dynamiques plausibles du point de vue économique.
    \item \textbf{Nature des sorties :} un GSE en risque neutre peut générer des trajectoires peu intuitives au regard de l'expérience historique, mais pertinentes pour la valorisation ; un GSE en monde réel doit au contraire produire des scénarios crédibles pour l'analyse prospective.
    \item \textbf{Usages en assurance vie :} le GSE en risque neutre est mobilisé pour le calcul du Best Estimate, de la TVOG et pour la valorisation des garanties, tandis que le GSE en monde réel est davantage utilisé dans le cadre de l'ORSA, de l'ALM et des projections de solvabilité à long terme.
\end{itemize}

Les deux approches sont donc complémentaires, dans la mesure où elles répondent à des objectifs distincts. Un assureur peut, par exemple, recourir à un GSE en risque neutre pour valoriser ses engagements, tout en utilisant un GSE en monde réel pour analyser l'évolution de ses risques et de sa solvabilité dans la durée. Il est toutefois essentiel de bien distinguer ces deux cadres, car leur confusion pourrait entraîner des erreurs significatives dans l'évaluation des engagements ou dans l'appréciation des risques.


\subsection{Vasicek}
\subsection{un deuxième modèle}
\subsection{le troisième modèle qui est le meilleur pour rendre compte de l'évolution des taux d'intérêt }

%%%%%%%%%%%%%%%%%%%%%%%%%%%%%%%%%
% Le modèle de gestion actif-passif
%%%%%%%%%%%%%%%%%%%%%%%%%%%%%%%%%
\section{Le modèle de gestion actif-passif}
\subsection{Les principes de l'ALM en assurance vie}
La gestion actif-passif ou Asset-Liability Management (ALM) est une approche stratégique utilisée par les institutions financières, notamment les compagnies d'assurance vie, pour gérer de manière intégrée leurs actifs et leurs passifs. L'objectif principal de l'ALM est d'assurer la solvabilité de l'institution tout en optimisant la rentabilité de son portefeuille d'investissement.
\subsection{La structure du modèle développé}
\subsubsection{La modélisation du passif}
\subsubsection{La modélisation de l'actif}

% \subsubsubsection{Les obligations}






% \subsubsubsection{Les actions}
% \subsubsubsection{L'immobilier}
\subsubsection{Le mécanisme d'allocation}
